\section{Two Sum}

Given an array of integers $nums$ and an integer $target$, find two indices $i$ and $j$ such that
$nums[i] + nums[j] = target$. You may assume that each input has exactly one solution, and you
may not use the same element twice.

\subsection{Input / Output}

\textbf{Input:} An array $nums$ of $n$ integers and an integer $target$.

\textbf{Output:} Two indices $i, j$ with $i \neq j$ such that $nums[i] + nums[j] = target$.

\subsection{Constraints}

$$2 \leq n \leq 10^4, \quad -10^9 \leq nums[i] \leq 10^9$$

\begin{diagram}[id="input-array", label="Input array visualization"]
\shape{nums}{Array}{size=5, data=[2,7,11,15,1], labels="0..4", label="$nums$"}
\end{diagram}

\begin{animation}[id="two-sum-brute", label="Two Sum -- Brute Force O(n^2)"]
\shape{data}{Array}{size=5, data=[2,7,11,15,1], labels="0..4", label="$nums$"}
\shape{status}{Array}{size=5, data=["","","","",""], labels="0..4", label="status"}
\compute{ target = 9 }

\step
\narrate{We want to find two numbers that sum to $target = 9$. We use the brute-force approach: try every pair $(i, j)$ with $i < j$.}

\step
\recolor{data.cell[0]}{state=current}
\recolor{status.cell[0]}{state=current}
\apply{status.cell[0]}{value="i=0"}
\annotate{data.cell[0]}{label="i", position=above, color=info}
\narrate{Outer loop: set $i = 0$. The value at index 0 is $nums[0] = 2$.}

\step
\recolor{data.cell[1]}{state=current}
\recolor{status.cell[1]}{state=current}
\apply{status.cell[1]}{value="j=1"}
\annotate{data.cell[1]}{label="j", position=above, color=info}
\annotate{data.cell[0]}{label="2+7=9?", position=below, color=warn}
\narrate{Inner loop: $j = 1$. Check pair $(0, 1)$: $nums[0] + nums[1] = 2 + 7 = 9$. That equals our target!}

\step
\recolor{data.cell[0]}{state=good}
\recolor{data.cell[1]}{state=good}
\recolor{status.cell[0]}{state=good}
\recolor{status.cell[1]}{state=good}
\apply{status.cell[0]}{value="found"}
\apply{status.cell[1]}{value="found"}
\annotate{data.cell[0]}{label="ans[0]", position=above, color=good}
\annotate{data.cell[1]}{label="ans[1]", position=above, color=good}
\annotate{data.cell[0]}{label="2+7=9", position=below, color=good}
\narrate{Match found! $nums[0] + nums[1] = 2 + 7 = 9 = target$. Return indices $[0, 1]$.}

\step
\recolor{data.cell[2]}{state=dim}
\recolor{data.cell[3]}{state=dim}
\recolor{data.cell[4]}{state=dim}
\recolor{status.cell[2]}{state=dim}
\recolor{status.cell[3]}{state=dim}
\recolor{status.cell[4]}{state=dim}
\apply{status.cell[2]}{value="skip"}
\apply{status.cell[3]}{value="skip"}
\apply{status.cell[4]}{value="skip"}
\narrate{The remaining elements are never checked because we already found our answer. In the worst case, though, we would continue the nested loop.}

\step
\narrate{Now let us trace what the \textbf{worst case} looks like. Consider $target = 16$ with the same array $[2, 7, 11, 15, 1]$.}

\step
\recolor{data.cell[0]}{state=current}
\recolor{data.cell[1]}{state=idle}
\recolor{data.cell[2]}{state=idle}
\recolor{data.cell[3]}{state=idle}
\recolor{data.cell[4]}{state=idle}
\recolor{status.cell[0]}{state=current}
\recolor{status.cell[1]}{state=idle}
\recolor{status.cell[2]}{state=idle}
\recolor{status.cell[3]}{state=idle}
\recolor{status.cell[4]}{state=idle}
\apply{status.cell[0]}{value="i=0"}
\apply{status.cell[1]}{value=""}
\apply{status.cell[2]}{value=""}
\apply{status.cell[3]}{value=""}
\apply{status.cell[4]}{value=""}
\annotate{data.cell[0]}{label="i", position=above, color=info}
\narrate{Worst-case trace with $target = 16$. Outer loop: $i = 0$. We must check $j = 1, 2, 3, 4$ before moving on.}

\step
\recolor{data.cell[1]}{state=current}
\apply{status.cell[1]}{value="2+7=9"}
\annotate{data.cell[1]}{label="j", position=above, color=warn}
\annotate{data.cell[1]}{label="9 != 16", position=below, color=error}
\narrate{$j=1$: $2 + 7 = 9 \neq 16$. No match. The cursor for $j$ advances.}

\step
\recolor{data.cell[1]}{state=dim}
\recolor{data.cell[2]}{state=current}
\apply{status.cell[1]}{value="x"}
\apply{status.cell[2]}{value="2+11=13"}
\annotate{data.cell[2]}{label="j", position=above, color=warn}
\annotate{data.cell[2]}{label="13 != 16", position=below, color=error}
\narrate{$j=2$: $2 + 11 = 13 \neq 16$. Still no match. Continue.}

\step
\recolor{data.cell[2]}{state=dim}
\recolor{data.cell[3]}{state=current}
\apply{status.cell[2]}{value="x"}
\apply{status.cell[3]}{value="2+15=17"}
\annotate{data.cell[3]}{label="j", position=above, color=warn}
\annotate{data.cell[3]}{label="17 != 16", position=below, color=error}
\narrate{$j=3$: $2 + 15 = 17 \neq 16$. Close but not equal. One more to check for $i = 0$.}

\end{animation}

\subsection{Complexity Analysis}

The brute-force algorithm uses two nested loops over the array:

$$T(n) = \sum_{i=0}^{n-2} \sum_{j=i+1}^{n-1} O(1) = \frac{n(n-1)}{2} = O(n^2)$$

\textbf{Time complexity:} $O(n^2)$ --- we examine every pair exactly once.

\textbf{Space complexity:} $O(1)$ --- no additional data structures are used.

A hash-map approach improves time to $O(n)$ at the cost of $O(n)$ space:

$$T_{hash}(n) = \sum_{i=0}^{n-1} O(1) = O(n)$$
